\section{讨论与总结}

\subsection{哪些改进得到本次结果支持}
这次实验并未得到一条“步骤更多就更好”的路线。单项比较表明，参考短段过滤丢失了较多仍可表达的原始范围；将点数与长度门槛分别检查后，两项调整都有作用。方向条件与过滤的组合曾在开发数据上取得额外收益，但新选择数据的三条几何退化使它退出，最终采用了更简单的S-D-P设置。

冻结后的600条记录支持这一选择，全量结果也显示覆盖增加明显大于存储量增加。主要贡献是降低参考过滤带来的信息损失，而不是发明新的DP算法。时间30秒、距离400工作米、最少2点、最短长度0、方向35度和DP5工作米，是在这次离散候选、数据与比较条件下选出的设置，不是对所有轨迹的最优参数承诺。

对LLM的使用也遵循同样的结果标准。建议会转成真实计算，合法性、局部误差和完整几何分开检查；允许反馈的模式据实测继续，而不是由模型自行宣布有效。记忆虽有实际送达与引用，现有比较未证明其因果质量收益。最后保留固定处理策略，并不否定工作流在组织实验、核验和修复中的作用，但不能据此宣称所有Agent组件都带来了可分离的效果提升。

\subsection{哪些问题尚不能由现有数据回答}
最大的限制是没有真实道路、运动类型和位置误差标签。共同原始参考能够帮助发现处理过程的信息损失，却不能区分所有噪声与真实运动；新增覆盖也可能包含较大偏差。方向反例和记录352的局部现象说明，单个算法的保证不能替代整条链的质量判断。

坐标基准及采样来源说明不完整，记录之间的真实个体独立性也不能确认。因此，本文结论限定于固定数学工作平面与所处理数据；分层样本结果、重复模型运行和全量生产分别解释，不相互替代。

实验采用有限参数域和少量有依据的结构调整，并未穷尽所有策略。它形成的主要认识是：\textbf{先让每项操作与评价可核对，再用共同条件下的实际结果决定取舍。}复杂候选可以被否定，原有设置也可以被调整；真正需要保留的是能够被数据支持的部分，而不是预先设想的方案形态。
